\emph{Pour chacune des affirmations suivantes, indiquer si elle est vraie ou fausse.}

	\medskip

	\emph{Chaque réponse doit être justifiée. Une réponse non justifiée ne rapporte aucun point.}

	\bigskip

	\begin{enumerate}[itemsep=8mm]
		\item On considère la fonction $f$ définie sur $]0 ~;~ +\infty [$ par : $f(x)=x \ln(x)$.

		\textbf{Affirmation 1 :}

\[\int_{1}^{\e} f(x) \,\mathrm{d} x=\dfrac{\e^{2}+1}{4}\]

\item Soient $n$ et $k$ deux entiers naturels non nuls tels que $k \leqslant n$.

\textbf{Affirmation 2 :}

\[n \times\binom{ n-1}{k-1}=k \times\binom{ n}{k}\]

		\item Pour les trois affirmations suivantes, on considère que l'espace est muni d'un repère orthonormé \Oijk.

\medskip

Soit $d$ la droite de représentation paramétrique : \quad $\left\{\begin{array}{l c l}
x&=&\phantom{2}t+1 \\
y&=&2 t+1 \\
z&=&-t
\end{array},\quad t \in \R\right.$.

		Soit $d'$ la droite de représentation paramétrique : \quad $\left\{\begin{array}{l c l}
x&=&2t'-1 \\
y&=&-t'+2 \\
z&=&\phantom{-}t'+1
\end{array},\quad t' \in \R\right.$.

\medskip

Soit $P$ le plan d'équation cartésienne : $2 x+y-2 z+18=0$.

Soit A le point de coordonnées $(-1 ~;~-3 ~;~ 2)$ et B le point de coordonnées $(-5 ~;~-5~;~6)$.

On appelle plan médiateur du segment [AB] le plan passant par le milieu du segment [AB] et orthogonal à la droite (AB).

\medskip

\textbf{Affirmation 3 :} Le point A appartient à la droite $d$.

\medskip

\textbf{Affirmation 4 :} Les droites $d$ et $d'$ sont sécantes.

\medskip

		\textbf{Affirmation 5 :} Le plan $P$ est le plan médiateur du segment [AB].
	\end{enumerate}
